% ============================================================
% Section: Application — Branch Banking Deregulation and
%          Income Inequality (Beck et al., 2010)
%
% Required packages: booktabs, threeparttable, amsmath,
%                    caption, multirow, geometry
% ============================================================

\section{Application: Branch Banking Deregulation and Income Inequality}
\label{sec:application}

We replicate and extend the application of \citet{beck2010} to illustrate
the sensitivity of difference-in-differences estimates to the choice of
estimator under staggered adoption with heterogeneous treatment effects.
\citet{beck2010} study how intrastate branch banking deregulation affected
income inequality across U.S.\ states from 1976 to 2006. Using a standard
two-way fixed effects (TWFE) regression, they find a negative and
statistically significant effect of deregulation on the log Gini
coefficient, concluding that branch banking deregulation reduced income
inequality. We show that this conclusion is fragile: once forbidden
comparisons are removed and the error covariance structure is accounted
for, the estimated effect is close to zero and statistically
indistinguishable from it.

\subsection{Data and Setting}

The sample comprises 49 U.S.\ states observed annually from 1976 to 2006,
yielding a balanced panel of 1{,}519 state-year observations. The outcome
variable is the log Gini coefficient (\texttt{ln\_gini}). Treatment
corresponds to the year in which a state permitted intrastate branch
banking (\texttt{branch\_reform}). Adoption is staggered: treatment cohorts
range from 1976 to 1999, and 12 states deregulated before the sample
period begins and are therefore ``always treated.''

We consider two samples. \textit{Panel~A} retains all 49 states, with
the always-treated states included and handled appropriately by each
estimator. \textit{Panel~B} restricts to the 36 states whose treatment
falls strictly within the sample period (1977--1999), excluding
always-treated states entirely. Panel~B provides a cleaner identification
setting because every state has at least one pre-treatment observation,
guaranteeing that individual counterfactual trends are identified for all
treated cohorts.

\subsection{Results}

Table~\ref{tab:beck_results} reports ATT estimates and standard errors
from eight estimators for both panels. The estimators fall into three
groups: (i) the standard TWFE regression; (ii) four heterogeneous-robust
estimators from the recent DiD literature---\citet{callaway2021} (CS),
\citet{sun2021} (SA), \citet{gardner2024} (Gardner), and the flexible TWFE
of \citet{wooldridge2025}; and (iii) our proposed GMM estimator under three
weighting schemes---identity (GMM-OLS), diagonal inverse-covariance
(GMM-Diag), and full inverse-covariance (GMM-Eff).

%-----------------------------------------------------------------
\begin{table}[htbp]
\centering
\caption{Effect of Branch Banking Deregulation on $\ln(\text{Gini})$}
\label{tab:beck_results}
\begin{threeparttable}
\begin{tabular}{lcccc}
\toprule
 & \multicolumn{2}{c}{\textit{Panel A} (49 states)}
 & \multicolumn{2}{c}{\textit{Panel B} (36 states)} \\
\cmidrule(lr){2-3}\cmidrule(lr){4-5}
Estimator & ATT & SE & ATT & SE \\
\midrule
\multicolumn{5}{l}{\textit{Standard estimator}} \\[2pt]
TWFE
  & $-0.0213^{***}$ & $(0.0076)$
  & $-0.0134^{*}$   & $(0.0071)$ \\[4pt]
\multicolumn{5}{l}{\textit{Heterogeneous-robust estimators}} \\[2pt]
Callaway \& Sant'Anna (2021)
  & $-0.0101$       & $(0.0078)$
  & $\phantom{-}0.0062$  & $(0.0053)$ \\
Sun \& Abraham (2021)
  & $-0.0354^{***}$ & $(0.0114)$
  & $-0.0536^{***}$ & $(0.0063)$ \\
Gardner et al.\ (2024)
  & $\phantom{-}0.0195^{***}$ & $(0.0067)$
  & $\phantom{-}0.0195^{***}$ & $(0.0067)$ \\
Flexible TWFE (Wooldridge 2025)
  & $-0.0165$       & $(0.0141)$
  & $\phantom{-}0.0195^{***}$ & $(0.0069)$ \\[4pt]
\multicolumn{5}{l}{\textit{Proposed GMM estimators (236 CATTs, 5{,}829 DiD moments)}} \\[2pt]
GMM-OLS (identity weights)
  & $\phantom{-}0.0266$       & $(0.0190)$
  & $\phantom{-}0.0266$       & $(0.0187)$ \\
GMM-Diag (diagonal $\Omega^{-1}$)
  & $\phantom{-}0.0247$       & $(0.0187)$
  & $\phantom{-}0.0247$       & $(0.0180)$ \\
GMM-Eff (full $\Omega^{-1}$)
  & $\phantom{-}0.0002$       & $(0.0141)$
  & $\phantom{-}0.0007$       & $(0.0140)$ \\
\bottomrule
\end{tabular}
\begin{tablenotes}[flushleft]\small
\item \textit{Notes:} Standard errors in parentheses.
  $^{*}p<0.10$,\ $^{**}p<0.05$,\ $^{***}p<0.01$ based on two-sided $z$-tests.
  TWFE, SA, CS, and Gardner standard errors are clustered by state.
  GMM-OLS standard errors use the sandwich formula
  $(Q'Q)^{-1}Q'\hat{\Omega}Q(Q'Q)^{-1}$;
  GMM-Diag uses $(Q'AQ)^{-1}Q'A\hat{\Omega}AQ(Q'AQ)^{-1}$
  with $A = \mathrm{diag}(\hat{\Omega})^{-1}$;
  GMM-Eff uses the delta-method formula $(Q'\hat{\Omega}^{-1}Q)^{-1}$.
  $\hat{\Omega}$ is estimated from autocovariances of
  treatment-adjusted, two-way-demeaned residuals.
  CS uses not-yet-treated states as the control group with doubly robust
  estimation and group-level aggregation.
  SA uses interaction-weighted aggregation with actual treatment cohorts
  for all states in both panels.
  Panel A GMM uses 37 eligible states (branch reform $\geq 1976$);
  Panel B GMM uses all 36 states.
\end{tablenotes}
\end{threeparttable}
\end{table}
%-----------------------------------------------------------------

\subsection{Why Estimates Differ Across Estimators}
\label{subsec:why_differ}

The wide dispersion of estimates across Table~\ref{tab:beck_results}
reflects three distinct sources of divergence: the bias from forbidden
comparisons embedded in TWFE, differences in control group definition and
aggregation across the heterogeneous-robust estimators, and differences
in how efficiently each estimator exploits the available moment conditions.

\paragraph{Forbidden comparisons and the TWFE bias.}
By the Goodman-Bacon decomposition \citep{goodman-bacon2021}, the TWFE
coefficient can be expressed as a variance-weighted average of all
possible $2\times 2$ difference-in-differences (DiD) estimates. In settings
with staggered adoption, this average includes ``forbidden comparisons''
in which an already-treated cohort serves as the control group. These
comparisons are biased because the control cohort's outcome trajectory
reflects its own accumulated treatment effect, not a clean counterfactual
for the focal cohort. In the Beck et al.\ application, the
Goodman-Bacon decomposition reveals that approximately 86 percent of the
total TWFE weight is assigned to forbidden comparisons involving
already-treated states as implicit controls \citep{arora2025}. The
strongly negative TWFE coefficient therefore does not reflect the true
average treatment effect; it reflects the direction of bias induced by
these contaminated comparisons.

\paragraph{Heterogeneous-robust estimators.}
The CS, Gardner, and Flexible TWFE estimators all eliminate forbidden
comparisons by construction, yet they yield markedly different results.
CS finds a negative but insignificant ATT in Panel A ($-0.0101$) and a
small positive one in Panel B ($+0.0062$), consistent with a near-zero
true effect. Gardner yields positive and significant estimates around
$+0.020$ in both panels. The Flexible TWFE estimator diverges sharply
across panels: it yields $-0.0165$ (insignificant) in Panel A but
$+0.0195$ (significant) in Panel B. This divergence arises because
the Flexible TWFE, unlike Gardner's two-stage approach, includes
always-treated states in the common time fixed-effects estimation.
Since these 12 states contribute only post-treatment observations,
they shift the estimated time trends and thereby alter the
counterfactual baseline for all other cohorts. Gardner avoids this
contamination by estimating unit and time fixed effects exclusively
from untreated observations in the first stage, so its estimate is
invariant to whether always-treated states are retained.

The SA estimator stands apart, yielding large, negative, and highly
significant estimates ($-0.0355$ in Panel A, $-0.0536$ in Panel B). This
divergence is attributable to its aggregation scheme. SA produces
interaction-weighted ATT estimates that average CATTs over event time
rather than over cohorts. If early-treatment cohorts have systematically
different underlying inequality trends---as is plausible for states that
deregulated before others---the event-time weighting can amplify cohort
composition effects. In Panel A, the inclusion of always-treated states
as an implicit never-treated reference group further distorts the baseline
comparison. The SA estimate should therefore be interpreted with caution
in this application.

The divergence between Gardner (positive, $\approx+0.020$) and CS
(near zero) reflects a subtler difference. Both eliminate forbidden
comparisons, but they exploit the data differently. Gardner relies on
the parallel trends assumption holding globally---it uses the
pre-treatment variation of all untreated units to estimate
counterfactual trends for all treated cohorts. CS estimates
cohort-specific counterfactuals using only the cohort's own
not-yet-treated comparisons. When treatment effect heterogeneity is
strong and cohort compositions differ, the two approaches can diverge.
Panel~B is the cleaner comparison here: in the absence of always-treated
contamination, Flexible TWFE agrees with Gardner ($+0.0195$), suggesting
that the imputation approach consistently delivers a positive estimate
once the data support it. The Panel~A discrepancy for Flexible TWFE
($-0.0165$) should be attributed to the distorted time fixed effects
caused by always-treated units rather than to a genuine treatment
effect estimate.

\paragraph{GMM weighting and efficiency.}
The three GMM estimators operate on an identical set of moment
conditions---5{,}829 not-yet-treated $2\times 2$ DiD comparisons spanning
236 cohort-time CATT cells---and differ only in the weighting matrix
applied to those moments.

GMM-OLS assigns equal weight to all 5{,}829 DiD moments, yielding
$\hat{\theta} \approx +0.027$. GMM-Diag upweights moments with low
marginal variance and downweights high-variance moments, yielding
$\hat{\theta} \approx +0.025$---a modest shift toward zero that reflects
the partial information in the diagonal structure of $\Omega$.
GMM-Eff applies the full inverse-covariance matrix $\Omega^{-1}$ as the
weighting matrix. This is the optimal weighting under
\citet{hansen1982}, and it accounts for the cross-moment correlations that
arise because the same state-year observation appears in multiple DiD
comparisons. The GMM-Eff estimate collapses to near zero
($\approx+0.0002$) in both panels.

This convergence toward zero as efficiency increases has a clear
interpretation. The 5{,}829 DiD moments are not independent: moments
sharing a focal cohort, a control cohort, or an overlapping time
window are positively correlated. Under equal weighting (GMM-OLS),
redundant moments receive excessive weight and pull the estimate away
from zero. The optimal weighting matrix $\Omega^{-1}$ effectively
orthogonalises the moment conditions, assigning near-zero weight to
moments that are nearly collinear with others and reserving
weight for the moments that carry orthogonal identifying information.
The resulting estimate reflects what the data actually identify after
this de-redundancy operation. The efficiency gain from full
inverse-covariance weighting over identity weighting is substantial:
GMM-Eff standard errors ($\approx 0.014$) are approximately 26 percent
smaller than GMM-OLS standard errors ($\approx 0.019$), consistent
with the simulation results in Section~\ref{sec:simulations}.

\subsection{Reliability of Estimates and Implications}
\label{subsec:reliability}

Not all estimates in Table~\ref{tab:beck_results} are equally reliable,
and it is important to distinguish between estimators on the basis of
their identifying assumptions and statistical properties.

\textit{TWFE is the least reliable.} The negative and significant TWFE
estimate is not credible as a causal effect estimate: its sign and
magnitude are driven overwhelmingly by forbidden comparisons, as the
Goodman-Bacon decomposition confirms. This is an instructive case study
in the mechanism by which staggered adoption, when combined with
heterogeneous effects, causes TWFE to recover a biased and potentially
sign-reversed estimate of the true ATT.

\textit{SA should be interpreted with caution.} The large negative SA
estimates appear inconsistent with the other heterogeneous-robust
estimators. As discussed above, the event-time aggregation weights
and, in Panel A, the treatment of always-treated states as
never-treated introduce plausible sources of distortion. Researchers
relying on SA in applications with always-treated units should verify
robustness of the reference group coding.

\textit{CS and Gardner are unbiased; Flexible TWFE is sensitive to
always-treated units.} CS and Gardner eliminate forbidden comparisons
and are consistent under parallel trends. The sign disagreement between
CS (near zero) and Gardner (positive) reflects genuine differences in
counterfactual trend construction: Gardner uses a global parallel
trends assumption while CS makes no cross-cohort extrapolation.
The CS estimate may therefore be the most conservative and credible
of the two. The Flexible TWFE estimate in Panel~A ($-0.0165$) should
be treated with caution: its Panel~A value is distorted by
always-treated states contaminating the time fixed effects, while its
Panel~B value ($+0.0195$) agrees with Gardner. In applications with
always-treated units, Flexible TWFE should either exclude those units
or be compared against a Panel~B counterpart for robustness.

\textit{GMM-Eff is the most reliable.} It is the minimum-variance
unbiased estimator among the GMM class under the parallel trends
assumption, exploiting all not-yet-treated comparisons while
accounting for the full error covariance structure. Its point
estimate of essentially zero ($+0.0002$ in Panel A, $+0.0007$ in
Panel B) is statistically indistinguishable from zero at any
conventional significance level, with the tightest standard errors
among the unbiased estimators. The GMM-Diag estimator, which uses
only the diagonal of $\Omega$, delivers an intermediate result:
it captures some but not all of the efficiency gains from accounting
for moment dependence, and its positive point estimate ($\approx +0.025$)
is a consequence of the residual weight assigned to
redundant, correlated DiD moments.

\medskip

\noindent
Taken together, the results in Table~\ref{tab:beck_results} deliver a
clear and consistent message from the unbiased estimators: once
forbidden comparisons are removed and the identifying information in
the data is extracted efficiently, there is no statistically
significant evidence that branch banking deregulation affected income
inequality. The original finding of \citet{beck2010} is an artifact
of the TWFE bias rather than a genuine causal effect. This
application illustrates that the choice of estimator is not merely a
technical detail in staggered DiD designs---it can determine whether
a paper's central conclusion is a positive, a negative, or a null
result.
